\section{Transition Function}

The transition function $\delta$ is defined as follows:
\begin{align*}
{%- for key, value in tf.items() -%}
    \delta({{ key[0] }},\,{{ key[1] }}) &= ({{ value[0] }},\,{{ value[1] }},\,{{ value[2] }}) {%- if not loop.last %}\\[0.5em]{%- endif %}
{%- endfor %}
\end{align*}

\section{Turing Machine Simulation Visualizations}

The simulation produced the following sequence of transitions. For each transition, the tape is displayed in a fixed window (so the tape does not shift), and the head (with its state label) moves along the tape. In each transition, we detail:
\begin{itemize}
    \item \textbf{Head Position:} The index of the cell where the head is located.
    \item \textbf{Read Symbol:} The symbol at that cell that is being processed.
    \item \textbf{Rewritten Symbol:} The symbol that is written on the tape (if it is changed).
    \item \textbf{New State:} The Turing machine's state after applying the transition.
\end{itemize}

{%- for idx in range(plots|length) %}
\subsection*{Transition {{ idx+1 }}}
\textbf{Explanation:}
{%- if idx == 0 %}
\newline In the initial configuration, the Turing Machine loads the input tape. The head is positioned at the starting cell (index \textbf{{ details[idx].head_position }}), reading the symbol \textbf{`{{ details[idx].read_symbol }}`}. Since no transition has been performed yet, no rewriting occurs.
{%- else %}
\newline In this transition:
\begin{itemize}
    \item The head is positioned at cell index \textbf{{ details[idx].head_position }}, reading the symbol \textbf{`{{ details[idx].read_symbol }}`}.
    \item The machine rewrites this cell with \textbf{`{{ details[idx].written_symbol }}`} (if different from the original symbol).
    \item The machine then moves the head and changes its state to \textbf{{ details[idx].new_state }}.
\end{itemize}
This step-by-step process is illustrated in the figure below.
{%- endif %}

\begin{figure}[h]
    \centering
    \includegraphics[width=0.8\linewidth]{plots/plot{{ idx+1 }}.png}
    \caption{Configuration at Transition {{ idx+1 }}}
\end{figure}

{%- endfor %}

\section{Halt and Final Tape}

At transition \(\mathbf{ {{ plots|length }} }\), the machine enters its halting state 
\(\mathbf{ {{ details[plots|length-1].new_state }} }\) and stops.

The final tape configuration is shown below:

\begin{figure}[h]
    \centering
    \includegraphics[width=0.8\linewidth]{plots/plot{{ plots|length }}.png}
    \caption{Final tape contents at halting state \(\mathbf{ {{ details[plots|length-1].new_state }} }\).}
\end{figure}
